\documentclass[a4paper,11pt]{article}

% ---------------------------------------------------------------------------
% Pacotes
% ---------------------------------------------------------------------------
\usepackage[utf8]{inputenc}
\usepackage[T1]{fontenc}
\usepackage[brazil]{babel}
\usepackage{geometry}
\usepackage{amsmath}
\usepackage{booktabs}
\usepackage{listings}
\usepackage{xcolor}
\usepackage{hyperref}
\usepackage{enumitem}

\geometry{margin=2.5cm}

\definecolor{codebg}{RGB}{245,245,245}
\definecolor{codekeyword}{RGB}{0,90,150}
\definecolor{codecomment}{RGB}{120,120,120}
\definecolor{codestring}{RGB}{150,60,20}

\lstdefinestyle{cpp}{
    language=C++,
    backgroundcolor=\color{codebg},
    basicstyle=\ttfamily\small,
    keywordstyle=\color{codekeyword}\bfseries,
    commentstyle=\color{codecomment}\itshape,
    stringstyle=\color{codestring},
    numbers=left,
    numberstyle=\tiny\color{gray},
    numbersep=6pt,
    breaklines=true,
    frame=single,
    tabsize=2,
    showstringspaces=false,
}
\lstset{style=cpp}

\hypersetup{
    colorlinks=true,
    linkcolor=blue,
    urlcolor=blue,
    citecolor=blue,
}

\title{Implementação de uma Biblioteca de Grafos em C++\\
\large Teoria dos Grafos (COS232) --- UFRJ}
\author{}
\date{}

\begin{document}
\maketitle

\begin{abstract}
Este documento descreve a implementação de uma biblioteca de grafos
não-direcionados em C++17, construída em torno de uma única interface
abstrata (\texttt{Graph}) com duas representações intercambiáveis ---
lista de adjacência e matriz de adjacência. O objetivo do projeto é
permitir comparar, de forma mensurável, o custo de cada representação
sobre as mesmas operações: construção, busca em largura e profundidade,
componentes conexas, distâncias e diâmetro (exato e aproximado).
\end{abstract}

\tableofcontents

% ===========================================================================
\section{Introdução}
% ===========================================================================

A biblioteca foi escrita para o trabalho de Teoria dos Grafos (COS232) e
não possui dependências externas. Todo o código público vive em um único
cabeçalho, \texttt{include/graph.hpp}, e cada função exportada documenta
argumentos, valor de retorno e custo assintótico \emph{para as duas
representações}, já que o mesmo algoritmo pode ter complexidades muito
diferentes dependendo de como o grafo está armazenado.

A organização dos arquivos é:

\begin{lstlisting}[language=bash,numbers=none]
include/graph.hpp     interface publica, com os docstrings
src/graph.cpp         AdjacencyList e AdjacencyMatrix
src/io.cpp            parser de entrada e escritores de CSV
src/search.cpp        bfs, dfs e os dois diametros
src/components.cpp    componentes conexas (largura e profundidade)
src/distance.cpp      getDistance
src/stats.cpp         computeStats
app/main.cpp          CLI que produz os CSVs de saida
app/benchmark.cpp      modos de tempo, memoria e relatorio
analysis/harness.py   roda o benchmark sobre config.csv
\end{lstlisting}

% ===========================================================================
\section{Arquitetura geral}
% ===========================================================================

O núcleo do design é a classe abstrata \texttt{Graph}, que define o
contrato que qualquer representação deve cumprir:

\begin{lstlisting}
class Graph {
public:
  int n; // numero de vertices
  int m; // numero de arestas

  explicit Graph(int vertexCount) : n(vertexCount), m(0) {}
  virtual ~Graph() = default;

  virtual void addEdge(int u, int v) = 0;
  virtual int degree(int u) const = 0;
  virtual std::vector<int> neighbors(int u) const = 0;
  virtual void finalize() = 0;
};
\end{lstlisting}

Os algoritmos (BFS, DFS, componentes, distâncias, diâmetro, estatísticas)
são todos funções livres que recebem \texttt{const Graph\&}, nunca o tipo
concreto. Isso significa que o mesmo código de \texttt{bfs} roda sobre
lista ou matriz sem alteração --- o que muda é apenas o custo de cada
chamada virtual, e é exatamente esse custo que o trabalho se propõe a
medir.

Vértices são indexados a partir de 1. Todo vetor interno é dimensionado
com tamanho $n+1$, deixando a posição 0 como preenchimento não utilizado;
um laço que começa em 0 lê lixo por engano, então essa convenção é
respeitada em toda a base de código.

% ===========================================================================
\section{Representações do grafo}
% ===========================================================================

\subsection{Lista de adjacência (\texttt{AdjacencyList})}

Guarda um \texttt{std::vector<std::vector<int>>}, um vetor de vizinhos por
vértice. Uma aresta repetida é armazenada novamente --- a estrutura se
comporta como um multigrafo --- e um laço adiciona 2 ao grau de \(u\),
preservando o lema do aperto de mãos ($\sum \deg(u) = 2m$) em relação ao
$m$ lido do arquivo. \texttt{finalize()} ordena cada linha para que os
vizinhos voltem em ordem ascendente, o que torna a saída das travessias
reprodutível.

\subsection{Matriz de adjacência (\texttt{AdjacencyMatrix})}

Guarda um \texttt{std::vector<std::vector<char>>}, um byte por par
ordenado de vértices. A escrita é idempotente: uma aresta repetida não
muda o grau, e um laço adiciona apenas 1 ao grau de \(u\). Por isso os
graus aqui descrevem o grafo simples subjacente, enquanto $m$ continua
contando o arquivo --- as duas representações discordam sempre que a
entrada repete uma aresta ou contém laços.

\subsection{Assimetria de custo entre representações}

Três detalhes de custo, fáceis de perder ao ler os pontos de chamada,
explicam por que a matriz é sistematicamente mais cara:

\begin{enumerate}
    \item \texttt{neighbors()} sempre retorna um \texttt{std::vector} novo
    por valor, alocando e copiando mesmo na lista, onde a linha já existe.
    \item \texttt{neighbors()} na matriz varre a linha inteira, $O(n)$,
    independentemente de quantos vizinhos $u$ realmente tem. Qualquer
    travessia que é $O(n+m)$ sobre lista vira $O(n^2)$ sobre matriz.
    \item \texttt{degree()} é $O(1)$ na lista mas $O(n)$ na matriz, então
    um laço sobre o grau de todos os vértices custa $O(n)$ contra a lista
    e $O(n^2)$ contra a matriz.
\end{enumerate}

\begin{table}[h]
\centering
\caption{Complexidade por operação, $n$ vértices, $m$ arestas, $k$ componentes.}
\begin{tabular}{@{}lll@{}}
\toprule
\textbf{Operação} & \textbf{Lista de adjacência} & \textbf{Matriz de adjacência} \\
\midrule
Construção            & $O(n)$              & $O(n^2)$ \\
\texttt{addEdge}       & $O(1)$ amortizado   & $O(1)$ \\
\texttt{degree(u)}     & $O(1)$              & $O(n)$ \\
\texttt{neighbors(u)}  & $O(\deg u)$         & $O(n)$ \\
\texttt{finalize}      & $O(m \log n)$       & $O(1)$ \\
\texttt{readGraph}     & $O(n + m \log n)$   & $O(n^2 + m)$ \\
\texttt{bfs} / \texttt{dfs} & $O(n + m)$    & $O(n^2)$ \\
\texttt{getDistance}   & $O(n + m)$          & $O(n^2)$ \\
Componentes conexas    & $O(n + m + k\log k)$ & $O(n^2)$ \\
\texttt{getApproximateDiameter} & $O(n+m)$  & $O(n^2)$ \\
\texttt{getExactDiameter} & $O(n \cdot (n+m))$ & $O(n^3)$ \\
Memória                & $O(n+m)$ inteiros   & $O(n^2)$ bytes \\
\bottomrule
\end{tabular}
\end{table}

Medido em um grafo de 10.000 vértices, um único BFS leva 584~\textmu s na
lista contra 41.327~\textmu s na matriz.

% ===========================================================================
\section{Estruturas de dados auxiliares}
% ===========================================================================

\begin{description}[style=nextline]
\item[\texttt{Edge}] par (\texttt{target}, \texttt{weight}), reservado
para suporte futuro a grafos com peso; não é usado pelo restante da
biblioteca hoje.
\item[\texttt{Representation}] \texttt{enum class} com os valores
\texttt{List} e \texttt{Matrix}, usado por \texttt{readGraph} para
decidir qual representação concreta instanciar.
\item[\texttt{SearchTree}] resultado de uma travessia: \texttt{parent} (0
para a raiz \emph{e} para vértices não alcançados) e \texttt{level}
($-1$ para vértices não alcançados). É \texttt{level}, não
\texttt{parent}, que distingue a raiz de um vértice fora do componente.
\item[\texttt{GraphStats}] agrega todas as estatísticas pedidas pelo
trabalho: contagens, grau mínimo/máximo/médio/mediano, número e tamanhos
extremos dos componentes, e os dois diâmetros.
\end{description}

% ===========================================================================
\section{Leitura e escrita}
% ===========================================================================

\subsection{Formato de entrada}

A primeira linha do arquivo é o número de vértices; cada linha seguinte é
uma aresta como um par de identificadores de vértice:

\begin{lstlisting}[language=bash,numbers=none]
5
1 2
2 5
5 3
4 5
1 5
\end{lstlisting}

Um inteiro solto ao final do arquivo é silenciosamente descartado.
\texttt{readGraph} lança \texttt{std::runtime\_error} se o arquivo não
abrir, se $n$ não puder ser lido, ou se $n < 0$.

\subsection{Saída}

Cada execução do CLI (\texttt{app/main.cpp}) escreve quatro CSVs em
\texttt{data/}, nomeados a partir do grafo de entrada:

\begin{table}[h]
\centering
\caption{Arquivos de saída gerados por execução.}
\begin{tabular}{@{}lll@{}}
\toprule
\textbf{Arquivo} & \textbf{Uma linha por} & \textbf{Colunas} \\
\midrule
\texttt{stats\_<nome>.csv} & representação & contagens, graus, componentes, diâmetros \\
\texttt{bfs\_tree\_<nome>.csv} & vértice & \texttt{node, parent, level} \\
\texttt{dfs\_tree\_<nome>.csv} & vértice & \texttt{node, parent, level} \\
\texttt{components\_<nome>.csv} & vértice & \texttt{node, componentId} \\
\bottomrule
\end{tabular}
\end{table}

% ===========================================================================
\section{Algoritmos de busca}
% ===========================================================================

\subsection{Busca em largura (\texttt{bfs})}

Percorre o grafo a partir de \texttt{source} e devolve uma
\texttt{SearchTree} em que \texttt{level[v]} é a distância real no grafo
entre a origem e $v$. Custo $O(n+m)$ na lista, $O(n^2)$ na matriz.

\subsection{Busca em profundidade (\texttt{dfs})}

Implementada de forma iterativa: a pilha guarda pares (vértice, pai), e um
vértice é assentado na árvore quando é desempilhado --- o que reproduz
exatamente a árvore de uma DFS recursiva. Aqui \texttt{level[v]} é a
\emph{profundidade de $v$ na árvore de DFS}, não a distância no grafo; para
distâncias deve-se usar \texttt{bfs} ou \texttt{getDistance}. O espaço
usado é $O(n+m)$, pois a pilha pode acumular até uma entrada por
extremidade de aresta, diferente do $O(n)$ que uma DFS recursiva usaria.

Existem ainda as variantes \texttt{bfsExploration} e
\texttt{dfsExploration}, que visitam todo vértice alcançável sem manter
estado --- não usadas atualmente pelo restante da biblioteca, mas
mantidas na interface. Vale notar que \texttt{dfsExploration} empilha
vizinhos em ordem ascendente enquanto \texttt{dfs} os empilha em ordem
descendente, de forma que as duas visitam um mesmo componente em ordens
diferentes.

% ===========================================================================
\section{Componentes conexas}
% ===========================================================================

\texttt{getComponentsByBreadth} e \texttt{getComponentsByDepth} particionam
o grafo rodando uma BFS (ou DFS) a partir de cada vértice ainda não
visitado. O resultado é um vetor de componentes, cada um em ordem de
descoberta, com os próprios componentes ordenados por tamanho
decrescente. Essa ordenação usa \texttt{std::sort} não estável, então
componentes de mesmo tamanho podem trocar de posição entre execuções ---
os identificadores de componente \textbf{não são estáveis} entre rodadas.

% ===========================================================================
\section{Distâncias e diâmetro}
% ===========================================================================

\subsection{Distância entre dois vértices}

\texttt{getDistance(g, u, v)} roda uma BFS completa a partir de \(u\) e
retorna o número de arestas no menor caminho até \(v\), ou $-1$ se $v$
não for alcançável.

\subsection{Diâmetro exato}

\texttt{getExactDiameter} roda uma BFS a partir de cada vértice e mantém o
maior nível finito encontrado --- a chamada mais cara da biblioteca:
$O(n \cdot (n+m))$ na lista, $O(n^3)$ na matriz. Em um grafo desconexo, o
valor textbook seria infinito; como pares inalcançáveis são ignorados,
o que se obtém é o maior diâmetro entre os componentes.

\subsection{Diâmetro aproximado}

\texttt{getApproximateDiameter} usa a técnica de \emph{double sweep}: uma
BFS a partir de um vértice qualquer, depois uma segunda BFS a partir do
vértice mais distante encontrado, reportando a excentricidade desse
segundo vértice. É um limite inferior do diâmetro --- exato em árvores,
subestimado em grafos gerais --- e varre apenas o maior componente, então
em um grafo desconexo pode ficar abaixo do diâmetro exato por um segundo
motivo: um componente menor pode conter o diâmetro real.

\subsection{Orçamento de trabalho}

Como \texttt{getExactDiameter} custa cerca de $n \cdot (n+m)$ inspeções de
aresta, \texttt{computeStats} recebe um orçamento nessas unidades
(\texttt{kExactDiameterBudget}, cerca de seis segundos em build \texttt{-O2}).
Acima do orçamento, \texttt{exactDiameter} volta como $-1$ e o grafo é
descrito apenas pela aproximação. Orçar o produto $n\cdot(n+m)$, em vez de
apenas $n$, é o que detecta grafos densos: dois grafos podem compartilhar
$n$ e ainda assim diferir em uma ordem de grandeza de custo.

% ===========================================================================
\section{Estatísticas (\texttt{computeStats})}
% ===========================================================================

Calcula todas as estatísticas pedidas pelo trabalho em uma única
passagem: contagens de vértices e arestas, grau mínimo, máximo, médio
($2m/n$, derivado de $m$ e não do vetor de graus) e mediano, número de
componentes com seus tamanhos extremos, e os dois diâmetros. O custo é
dominado pelo diâmetro exato quando ele roda; quando é pulado, o laço de
graus e a ordenação dominam ($O(n\log n + m)$ na lista, $O(n^2)$ na
matriz).

% ===========================================================================
\section{Aplicações: CLI e benchmark}
% ===========================================================================

\subsection{\texttt{app/main.cpp}}

Ponto de entrada que produz os quatro CSVs de saída. Sem argumentos, lê
\texttt{config.csv} linha a linha (\texttt{grafo,representação}),
resolvendo \texttt{<nome>} para \texttt{data/<nome>.txt}; com argumentos,
processa apenas o grafo indicado. Um grafo que falha é reportado e
pulado, sem abortar o restante da execução.

\subsection{\texttt{app/benchmark.cpp}}

Produz as medições de desempenho em quatro modos:

\begin{itemize}
    \item \texttt{time} --- média e mediana de 100 execuções de BFS e de
    DFS, a partir da mesma amostra aleatória de vértices, medindo só o
    algoritmo.
    \item \texttt{mem} --- mantém o grafo carregado e sinaliza
    \texttt{READY} para que um processo externo capture o RSS.
    \item \texttt{report} --- respostas específicas do estudo de caso:
    pais em árvores de BFS/DFS, distâncias entre pares, contagem e
    tamanhos de componentes, e os dois diâmetros.
    \item \texttt{all} --- os três anteriores em um único processo, com
    a medição de memória primeiro, antes que tempo e relatório aloquem
    suas próprias estruturas.
\end{itemize}

\texttt{analysis/harness.py} conduz o \texttt{benchmark} sobre todo o
\texttt{config.csv} e agrega os resultados com \texttt{pandas}.

% ===========================================================================
\section{Convenções de projeto}
% ===========================================================================

\begin{itemize}
    \item \textbf{$-1$ sempre significa ``indefinido''}: o \texttt{level}
    de um vértice não alcançado, uma \texttt{getDistance} entre
    componentes diferentes, um \texttt{exactDiameter} pulado por
    orçamento.
    \item \textbf{O \texttt{level} de uma DFS é profundidade na árvore},
    não distância no grafo.
    \item \textbf{\texttt{parent} vale 0 tanto para a raiz quanto para
    vértices inalcançáveis} --- só \texttt{level} distingue os dois
    casos.
    \item \textbf{Nenhuma das duas representações deduplica arestas.} A
    lista guarda multiplicidade; a escrita da matriz é idempotente. As
    duas discordam em grau sempre que a entrada repete uma aresta ou
    contém laços.
    \item \textbf{\texttt{getApproximateDiameter} é um limite inferior},
    exato em árvores, subestimado em grafos gerais.
    \item \textbf{Identificadores de componente não são estáveis entre
    execuções}, pois a ordenação por tamanho usa \texttt{std::sort}, que
    não é estável.
\end{itemize}

% ===========================================================================
\section{Diagrama de classes}
% ===========================================================================

A Figura~\ref{fig:classes} (fora deste documento, em
\texttt{docs/diagrama\_classes.mmd}) descreve a estrutura estática do
código em notação Mermaid. Resumidamente:

\begin{itemize}
    \item \texttt{Graph} é a interface abstrata; \texttt{AdjacencyList} e
    \texttt{AdjacencyMatrix} a implementam.
    \item \texttt{Edge}, \texttt{SearchTree} e \texttt{GraphStats} são
    estruturas de dados simples (\emph{value types}), sem comportamento.
    \item Os algoritmos (\texttt{bfs}, \texttt{dfs}, componentes,
    distâncias, diâmetros, \texttt{computeStats}, leitura/escrita de
    CSV) são funções livres que operam sobre \texttt{const Graph\&} ---
    não pertencem a nenhuma classe, por isso aparecem agrupadas no
    diagrama como um módulo \texttt{GraphAlgorithms} apenas para fins de
    visualização, não como uma classe real do código.
\end{itemize}

O arquivo-fonte do diagrama pode ser renderizado em
\url{https://mermaid.live} ou por qualquer visualizador Mermaid
compatível (GitHub, VS Code, etc.):

\begin{lstlisting}[language=bash,numbers=none,caption={Trecho do diagrama de classes (ver docs/diagrama\_classes.mmd para o arquivo completo)}]
classDiagram
    class Graph {
        <<abstract>>
        +int n
        +int m
        +addEdge(int u, int v) void
        +degree(int u) int
        +neighbors(int u) vector~int~
        +finalize() void
    }
    Graph <|-- AdjacencyList
    Graph <|-- AdjacencyMatrix
\end{lstlisting}

% ===========================================================================
\section{Conclusão}
% ===========================================================================

A biblioteca separa deliberadamente \emph{interface} (o contrato
\texttt{Graph}) de \emph{implementação} (lista vs. matriz), e separa
\emph{estrutura de dados} de \emph{algoritmo}, mantendo BFS, DFS,
componentes, distâncias e diâmetro como funções livres independentes da
representação concreta. Essa escolha de design é o que torna possível
medir, com o mesmo código de algoritmo, o custo real de cada
representação --- o objetivo central do trabalho.

\end{document}
